\documentclass[12pt,a4paper]{article}
\usepackage[utf8]{inputenc}
\usepackage[T1]{fontenc}
\usepackage{amsmath}
\usepackage{amssymb}
\usepackage{graphicx}
\usepackage[french]{babel}
\usepackage[a4paper, margin=2.5cm]{geometry}
\usepackage[european]{circuitikz}
\usetikzlibrary{babel}
\usepackage{siunitx} % Pour les unités SI
\usepackage{pgfplots}
\pgfplotsset{compat=1.17}

\title{Corrigé de l'Examen d'Électronique : Diagramme de Bode}
\author{Gemini Code Assist}
\date{\today}

\begin{document}

\maketitle

\section{Analyse qualitative (Sans calcul !)}

\subsection*{a. Cas où \(\omega \to 0\)}
À très basse fréquence (en régime continu), le condensateur \(C\) se comporte comme un circuit ouvert (\(Z_C = \frac{1}{j\omega C} \to \infty\)).

\begin{center}
\begin{circuitikz}[scale=1.2, transform shape]
    \draw (0,2) node[circ, label=left:P] (P) {};
    \draw (0,0) node[circ, label=left:Q] (Q) {} node[ground]{};
    \node at (P) [above right=0.1cm] {\(V_1\)};
    \draw (P) to[R, l=$R$] (3,2) node[circ, label=above:N]{};
    \draw (3,2) to[R, l=$R$] (3,0) node[ground]{};
    \draw (3,2) to[C, l=$C$] (5,2) node[circ, label=above:A]{};
    \draw (5,2) to[R, l=$R_L$] (5,0) node[circ, label=right:B] (B) {} node[ground]{};
    \draw[->] (5.8,0) -- (5.8,2) node[midway, right] {$V_2$};
\end{circuitikz}
\end{center}

Aucun courant ne peut circuler à travers le condensateur \(C\) pour atteindre le nœud A. Par conséquent, aucun courant ne traverse la résistance de charge \(R_L\). La tension à ses bornes est donc nulle.
\begin{equation}
    V_2 = 0 \, \text{V}
\end{equation}

\subsection*{b. Cas où \(\omega \to \infty\)}
À très haute fréquence, le condensateur \(C\) se comporte comme un court-circuit (\(Z_C = \frac{1}{j\omega C} \to 0\)). Le nœud central (que nous appellerons N) est donc directement connecté au nœud A.

\begin{center}
\begin{circuitikz}[scale=1.2, transform shape]
    \draw (0,2) node[circ, label=left:P] (P) {};
    \draw (0,0) node[circ, label=left:Q] (Q) {} node[ground]{};
    \node at (P) [above right=0.1cm] {\(V_1\)};
    \draw (P) to[R, l=$R$] (3,2) node[circ, label=above:N]{};
    \draw (3,2) to[R, l=$R$] (3,0) node[ground]{};
    \draw (3,2) to[short] (5,2) node[circ, label=above:A]{};
    \draw (5,2) to[R, l=$R_L$] (5,0) node[circ, label=right:B] {} node[ground]{};
    \draw[->] (5.8,0) -- (5.8,2) node[midway, right] {$V_2$};
\end{circuitikz}
\end{center}

Le circuit devient un diviseur de tension. La tension \(V_2\) (qui est la tension au nœud A) est donnée par la résistance \(R\) en parallèle avec \(R_L\), divisée par la somme de la première résistance \(R\) et de cette combinaison parallèle.
\begin{equation}
    V_2 = V_1 \cdot \frac{R \parallel R_L}{R + (R \parallel R_L)} = V_1 \cdot \frac{\frac{R \cdot R_L}{R+R_L}}{R + \frac{R \cdot R_L}{R+R_L}} = V_1 \cdot \frac{R \cdot R_L}{R(R+R_L) + R \cdot R_L}
\end{equation}
\begin{equation}
    \frac{V_2}{V_1} = \frac{R \cdot R_L}{R^2 + R \cdot R_L + R \cdot R_L} = \frac{R_L}{R + 2R_L}
\end{equation}

\subsection*{c. Quel type de montage est-ce ?}
Le gain est nul à basse fréquence et tend vers une valeur constante non nulle (\(\frac{R_L}{R + 2R_L}\)) à haute fréquence. Il s'agit donc d'un \textbf{filtre passe-haut} du premier ordre.

\section{En fréquentiel (Impédances complexes)}

\subsection*{a. Appliquer le théorème de Thévenin entre A et B}
On cherche l'équivalent de Thévenin du circuit vu depuis la charge \(R_L\), c'est-à-dire entre les points A et B (B étant la masse).

\paragraph{Tension de Thévenin (\(E_T\)):}
C'est la tension à vide \(V_{AB}\) lorsque \(R_L\) est déconnectée.

\begin{center}
\begin{circuitikz}[scale=1.2, transform shape]
    \draw (0,2) node[circ, label=left:P] (P) {};
    \draw (0,0) node[circ, label=left:Q] (Q) {} node[ground]{};
    \node at (P) [above right=0.1cm] {\(V_1\)};
    \draw (P) to[R, l=$R$] (3,2) node[circ, label=above:N]{};
    \draw (3,2) to[R, l=$R$] (3,0) node[ground]{};
    \draw (3,2) to[C, l=$C$] (5,2) node[circ, label=above:A]{};
    \node at (6.3, 2.3) {\(E_T = V_{A, \text{vide}}\)};
\end{circuitikz}
\end{center}

À vide, aucun courant ne circule dans \(C\). La tension en A est donc la même que la tension au nœud N. Le nœud N est au milieu d'un pont diviseur de tension formé par les deux résistances \(R\).
\begin{equation}
    E_T = V_{N} = V_1 \cdot \frac{R}{R+R} = \frac{V_1}{2}
\end{equation}

\paragraph{Impédance de Thévenin (\(Z_T\)):}
On 'éteint' la source de tension \(V_1\) (on la remplace par un court-circuit) et on calcule l'impédance vue depuis A.

\begin{center}
\begin{circuitikz}[scale=1.2, transform shape]
    \draw (0,2) node[circ, label=left:P] (P) {} node[ground]{};
    \draw (P) to[R, l=$R$] (3,2) node[circ, label=above:N]{};
    \draw (3,2) to[R, l=$R$] (3,0) node[ground]{};
    \draw (3,2) to[C, l=$C$] (5,2) node[circ, label=above:A]{};
    \draw[->, thick] (6.5,2) -- (5.5,2);
    \node at (7,2) {\(Z_T\)};
\end{circuitikz}
\end{center}

Vue depuis A, on voit le condensateur \(C\) en série avec l'impédance du nœud N vers la masse. Depuis N, les deux résistances \(R\) sont en parallèle (l'une vers P qui est à la masse, l'autre vers Q qui est à la masse).
\begin{equation}
    Z_T = Z_C + (R \parallel R) = \frac{1}{j\omega C} + \frac{R}{2}
\end{equation}

\subsection*{b. Nouveau schéma équivalent}
Le circuit complet peut être remplacé par le modèle de Thévenin suivi de la charge \(R_L\).

\begin{center}
\begin{circuitikz}[scale=1.2, transform shape]
    \draw (0,0) node[ground]{} to[sV, v<=$E_T$] (0,2) to[generic, l=$Z_T$] (3,2) to[R, l=$R_L$] (3,0) node[ground]{};
    \draw[<->] (3.3, 2) -- (3.3, 0);
    \node at (3.7, 1) {\(V_2\)};
\end{circuitikz}
\end{center}

\subsection*{c. Déduire la relation entre V1 et V2}
En utilisant le schéma équivalent, \(V_2\) est obtenue par un simple diviseur de tension.
\begin{equation}
    V_2 = E_T \cdot \frac{R_L}{Z_T + R_L}
\end{equation}
En substituant les expressions de \(E_T\) et \(Z_T\) :
\begin{equation}
    V_2 = \frac{V_1}{2} \cdot \frac{R_L}{\frac{R}{2} + \frac{1}{j\omega C} + R_L}
\end{equation}
La fonction de transfert \(H(j\omega) = V_2/V_1\) est donc :
\begin{equation}
    H(j\omega) = \frac{R_L/2}{R_L + R/2 + \frac{1}{j\omega C}} = \frac{j\omega C (R_L/2)}{1 + j\omega C (R_L + R/2)}
\end{equation}
Cette expression confirme bien que le circuit est un filtre passe-haut du premier ordre.

\section{Fonction de transfert}

\textbf{Note importante :} Il existe une incohérence entre le circuit de la Figure 1 et la forme de la fonction de transfert proposée à la question 3. L'analyse précédente (questions 1 et 2) montre que le circuit est un \textit{filtre passe-haut}. Cependant, la fonction \(H(j\omega)\) proposée ci-dessous est celle d'un \textit{filtre passe-tout} (ou déphaseur). De plus, le gain statique \(G_0\) est négatif, ce qui est physiquement impossible pour ce circuit passif.

\textbf{Pour la suite de l'exercice, nous supposerons que la fonction de transfert donnée est celle qu'il faut étudier, indépendamment du schéma initial.}

\subsection*{a. Calculer la fonction de transfert \(H(j\omega) = V_2/V_1\)}
On nous demande d'utiliser la forme :
\begin{equation}
    H(j\omega) = G_0 \frac{1 - j\frac{\omega}{\omega_0}}{1 + j\frac{\omega}{\omega_0}}
\end{equation}

\subsection*{b. Montrer que \(G_0 = -\frac{R_L}{2R_L+R}\)}
Cette expression est donnée dans le cadre de la définition du problème pour la fonction passe-tout. Elle ne peut pas être démontrée à partir du circuit passe-haut de la Figure 1. Nous l'acceptons comme une donnée.

\subsection*{c. Calcul de R pour \(f_0 = \SI{380}{\hertz}\)}
On nous donne \(\omega_0 = \frac{1}{(R_L + R/2)C}\), \(C = \SI{15}{\nano\farad}\), \(R_L = \SI{22}{\kilo\ohm}\) et \(f_0 = \SI{380}{\hertz}\).

On sait que \(\omega_0 = 2\pi f_0\). On peut donc isoler R :
\begin{align*}
    \omega_0 &= \frac{1}{(R_L + R/2)C} \\
    R_L + R/2 &= \frac{1}{\omega_0 C} = \frac{1}{2\pi f_0 C} \\
    R/2 &= \frac{1}{2\pi f_0 C} - R_L \\
    R &= 2 \left( \frac{1}{2\pi f_0 C} - R_L \right)
\end{align*}
Application numérique :
\begin{align*}
    R &= 2 \left( \frac{1}{2\pi \cdot 380 \cdot 15 \times 10^{-9}} - 22000 \right) \\
    R &= 2 \left( \frac{1}{3.5814 \times 10^{-5}} - 22000 \right) \\
    R &= 2 \left( 27922.5 - 22000 \right) \\
    R &= 2 \left( 5922.5 \right) \\
    R &= 11845 \, \Omega \approx \SI{11.85}{\kilo\ohm}
\end{align*}

\section{Tracé du Diagramme de Bode}
On considère \(H(j\omega)\) avec \(f_0 = \SI{380}{\hertz}\) et \(G_0 = 0.4\).
La fonction est donc \(H(j\omega) = 0.4 \frac{1 - j\frac{f}{f_0}}{1 + j\frac{f}{f_0}}\).

\subsection*{a. Expression du module et de l'argument}
\paragraph{Module :}
Le module de la fonction de transfert est :
\begin{equation}
    |H(j\omega)| = |0.4| \cdot \frac{|1 - j\frac{f}{f_0}|}{|1 + j\frac{f}{f_0}|} = 0.4 \cdot \frac{\sqrt{1^2 + (-\frac{f}{f_0})^2}}{\sqrt{1^2 + (\frac{f}{f_0})^2}} = 0.4
\end{equation}
Le module est constant et ne dépend pas de la fréquence.
En décibels, le gain est :
\begin{equation}
    H_{dB} = 20 \log_{10}(|H(j\omega)|) = 20 \log_{10}(0.4) \approx \SI{-7.96}{\deci\bel}
\end{equation}

\paragraph{Argument (\(\phi\)):}
L'argument (ou la phase) de la fonction de transfert est :
\begin{align*}
    \phi(\omega) &= \arg(0.4) + \arg\left(1 - j\frac{\omega}{\omega_0}\right) - \arg\left(1 + j\frac{\omega}{\omega_0}\right) \\
    \phi(\omega) &= 0 + \left(-\arctan\left(\frac{\omega}{\omega_0}\right)\right) - \left(\arctan\left(\frac{\omega}{\omega_0}\right)\right) \\
    \phi(\omega) &= -2 \arctan\left(\frac{\omega}{\omega_0}\right) = -2 \arctan\left(\frac{f}{f_0}\right)
\end{align*}

\subsection*{b. Tracer le diagramme de bode du module}
Le gain en module \(|H(j\omega)|\) est constant et égal à 0.4. Le diagramme de gain est donc une ligne horizontale à \(\approx \SI{-7.96}{dB}\) sur toute la plage de fréquences de \SI{10}{\hertz} à \SI{10}{\kilo\hertz}.

\begin{center}
    \begin{tikzpicture}
        \begin{semilogxaxis}[
            width=0.8\textwidth,
            height=6cm,
            xlabel={Fréquence $f$ (Hz)},
            ylabel={Gain $H_{dB}$ (dB)},
            grid=both,
            xmin=10, xmax=10000
        ]
        \addplot[blue, thick, domain=10:10000] {20*log10(0.4)};
        \end{semilogxaxis}
    \end{tikzpicture}
    \textit{Le diagramme de gain est une droite horizontale à -7.96 dB.}
\end{center}

\subsection*{c. Tracer le diagramme de bode de la phase}
La phase est \(\phi(f) = -2 \arctan(f/380)\).
\begin{itemize}
    \item Pour \(f \ll f_0\) (ex: \SI{10}{\hertz}), \(\phi \approx 0^\circ\).
    \item Pour \(f = f_0 = \SI{380}{\hertz}\), \(\phi = -2 \arctan(1) = -2 \cdot 45^\circ = -90^\circ\).
    \item Pour \(f \gg f_0\) (ex: \SI{10}{\kilo\hertz}), \(\phi \approx -2 \cdot 90^\circ = -180^\circ\).
\end{itemize}

Calcul de quelques points (avec \(f_0 = 380\)):
\begin{center}
    \resizebox{\textwidth}{!}{%
    \begin{tabular}{|l|c|c|c|c|c|c|}
    \hline
    \textbf{f (Hz)} & 30 & 100 & \textbf{380} (\(f_0\)) & 1000 & 3000 & 10000 \\
    \hline
    \(f/f_0\) & 0.08 & 0.26 & 1 & 2.63 & 7.89 & 26.3 \\
    \hline
    \(\phi(f)\) (deg) & $-9.0^\circ$ & $-29.4^\circ$ & \textbf{-90$^\circ$} & $-138.4^\circ$ & $-165.6^\circ$ & $-175.6^\circ$ \\
    \hline
    \end{tabular}%
    }
\end{center}

\begin{center}
    \begin{tikzpicture}
        \begin{semilogxaxis}[
            width=0.8\textwidth,
            height=6cm,
            xlabel={Fréquence $f$ (Hz)},
            ylabel={Phase $\phi$ (degrés)},
            grid=both,
            xmin=10, xmax=10000,
            ytick={0, -45, -90, -135, -180}
        ]
        \addplot[red, thick, domain=10:10000, samples=200] {-2*atan(x/380)};
        \end{semilogxaxis}
    \end{tikzpicture}
    \textit{Le diagramme de phase part de 0°, passe par -90° à \(f_0\), et tend vers -180°.}
\end{center}

\subsection*{d. Commenter le comportement en fréquence du montage}
En se basant sur l'étude de la fonction de transfert des questions 3 et 4, le montage se comporte comme un \textbf{filtre passe-tout} du premier ordre (aussi appelé déphaseur pur).

Son gain en amplitude est constant sur toute la plage de fréquence, ce qui signifie qu'il ne filtre (n'atténue) aucune fréquence. Son seul effet est d'introduire un déphasage qui varie avec la fréquence, allant de 0° à -180°.

Ce comportement est en \textbf{contradiction directe} avec la réponse à la question 1c. L'analyse du circuit initial avait montré un comportement de filtre passe-haut. Cela confirme l'incohérence entre le schéma et les équations fournies dans l'énoncé de l'examen.

\end{document}